% GENERATED FILE. Edit canonical lessons, then run npm run generate:books.
% canonical-source-hash: fnv1a64:290e554fc93249ac
% canonical-lessons: HI-C532-git, HI-C532-lay, HI-C532-riyaz, HI-C532-mahfil, HI-C532-nacana

\chapter{A Song to Dance To}
\label{ch:hi-a-song-to-dance-to}
\hlchaptermodality{532}

\begin{chapteropening}
\textbf{By the end of this chapter:} \emph{I can talk about songs, rhythm and dancing.}

Complete the last lesson of chapter 532: I can talk about songs, rhythm and dancing.
\end{chapteropening}

\section[\texorpdfstring{\hi{गीत} (gīt)}{गीत (gīt)}]{\hi{गीत} (gīt) — a song, a lyric}

\label{lesson:HI-C532-git}



\begin{quote}
Before the new one: say the Hindi for a musical instrument, then the Hindi for a mridang (a barrel drum).
\end{quote}

\subsection*{You'll want to know: \hi{गीत}}
\textbf{\hi{गीत}} — \emph{gīt} — \textquotedblleft{}a song, a lyric\textquotedblright{}.

\begin{quote}
\hi{शादी} \hi{में} \hi{औरतों} \hi{ने} \hi{पुराने} \hi{गीत} \hi{गाए।}
\end{quote}

\begin{quote}
\emph{Śādī meṁ aurtoṁ ne purāne gīt gāe.}
\end{quote}

\begin{quote}
At the wedding the women sang old songs.
\end{quote}

Masculine.

\begin{cousinweb}
\textbf{\hi{गीत}} (\emph{gīt}) is a tatsama, taken whole from Sanskrit \emph{gīta}, \textquotedblleft{}sung; a song\textquotedblright{} (Monier-Williams).
\end{cousinweb}

\begin{grammarlens}[title={What you've built}]
One more word for describing what you lived and felt.
\end{grammarlens}

\subsection*{Your turn}
\begin{itemize}
\raggedright
  \item \emph{Say it:} \emph{gīt}
  \item \emph{Say it:} \emph{gīt}, once more
  \item \emph{Say it:} \emph{Śādī meṁ aurtoṁ ne purāne gīt gāe.} — At the wedding the women sang old songs.
  \item \emph{Recall:} say the Hindi for a musical instrument, then the Hindi for a mridang (a barrel drum), then say \emph{gīt} again
\end{itemize}

\begin{tcolorbox}[breakable,colback=teal!4,colframe=teal!35!black,title={Before you move on}]
What does \hi{गीत} mean? (\textquotedblleft{}A song, a lyric\textquotedblright{}.) Say it once more.
\end{tcolorbox}
\section[\texorpdfstring{\hi{लय} (lay)}{लय (lay)}]{\hi{लय} (lay) — tempo; rhythm}

\label{lesson:HI-C532-lay}



\begin{quote}
Before the new one: say the Hindi for a mridang (a barrel drum), then the Hindi for a lyric.
\end{quote}

\subsection*{You'll want to know: \hi{लय}}
\textbf{\hi{लय}} — \emph{lay} — \textquotedblleft{}tempo; rhythm\textquotedblright{}.

\begin{quote}
\hi{गाने} \hi{की} \hi{लय} \hi{बहुत} \hi{तेज़} \hi{है।}
\end{quote}

\begin{quote}
\emph{Gāne kī lay bahut tez hai.}
\end{quote}

\begin{quote}
The song's tempo is very fast.
\end{quote}

Feminine.

\begin{cousinweb}
\textbf{\hi{लय}} (\emph{lay}) is a tatsama, taken whole from Sanskrit \emph{laya}, \textquotedblleft{}(in music) time\textquotedblright{} (Monier-Williams).
\end{cousinweb}

\begin{grammarlens}[title={What you've built}]
One more word for describing what you lived and felt.
\end{grammarlens}

\subsection*{Your turn}
\begin{itemize}
\raggedright
  \item \emph{Say it:} \emph{lay}
  \item \emph{Say it:} \emph{lay}, once more
  \item \emph{Say it:} \emph{Gāne kī lay bahut tez hai.} — The song's tempo is very fast.
  \item \emph{Recall:} say the Hindi for a mridang (a barrel drum), then the Hindi for a lyric, then say \emph{lay} again
\end{itemize}

\begin{tcolorbox}[breakable,colback=teal!4,colframe=teal!35!black,title={Before you move on}]
What does \hi{लय} mean? (\textquotedblleft{}Tempo; rhythm\textquotedblright{}.) Say it once more.
\end{tcolorbox}
\section[\texorpdfstring{\hi{रियाज़} (riyāz)}{रियाज़ (riyāz)}]{\hi{रियाज़} (riyāz) — practice (of music)}

\label{lesson:HI-C532-riyaz}



\begin{quote}
Before the new one: say the Hindi for a lyric, then the Hindi for tempo.
\end{quote}

\subsection*{You'll want to know: \hi{रियाज़}}
\textbf{\hi{रियाज़}} — \emph{riyāz} — \textquotedblleft{}practice (of music)\textquotedblright{}.

\begin{quote}
\hi{अच्छा} \hi{गायक} \hi{बनने} \hi{के} \hi{लिए} \hi{रोज़} \hi{रियाज़} \hi{करना} \hi{पड़ता} \hi{है।}
\end{quote}

\begin{quote}
\emph{Acchā gāyak banne ke liye roz riyāz karnā paṛtā hai.}
\end{quote}

\begin{quote}
To become a good singer you have to practise every day.
\end{quote}

Masculine.

\begin{cousinweb}
\textbf{\hi{रियाज़}} (\emph{riyāz}) is usually traced to Arabic \emph{riyāḍa}, \textquotedblleft{}exercise, training\textquotedblright{}, through Persian.
\end{cousinweb}

\begin{grammarlens}[title={What you've built}]
One more word for describing what you lived and felt.
\end{grammarlens}

\subsection*{Your turn}
\begin{itemize}
\raggedright
  \item \emph{Say it:} \emph{riyāz}
  \item \emph{Say it:} \emph{riyāz}, once more
  \item \emph{Say it:} \emph{Acchā gāyak banne ke liye roz riyāz karnā paṛtā hai.} — To become a good singer you have to practise every day.
  \item \emph{Recall:} say the Hindi for a lyric, then the Hindi for tempo, then say \emph{riyāz} again
\end{itemize}

\begin{tcolorbox}[breakable,colback=teal!4,colframe=teal!35!black,title={Before you move on}]
What does \hi{रियाज़} mean? (\textquotedblleft{}Practice (of music)\textquotedblright{}.) Say it once more.
\end{tcolorbox}
\section[\texorpdfstring{\hi{महफ़िल} (mahfil)}{महफ़िल (mahfil)}]{\hi{महफ़िल} (mahfil) — a gathering (for music or poetry)}

\label{lesson:HI-C532-mahfil}



\begin{quote}
Before the new one: say the Hindi for tempo, then the Hindi for practice (of music).
\end{quote}

\subsection*{You'll want to know: \hi{महफ़िल}}
\textbf{\hi{महफ़िल}} — \emph{mahfil} — \textquotedblleft{}a gathering (for music or poetry)\textquotedblright{}.

\begin{quote}
\hi{शाम} \hi{को} \hi{घर} \hi{पर} \hi{संगीत} \hi{की} \hi{महफ़िल} \hi{थी।}
\end{quote}

\begin{quote}
\emph{Śām ko ghar par saṅgīt kī mahfil thī.}
\end{quote}

\begin{quote}
In the evening there was a musical gathering at home.
\end{quote}

Feminine.

\begin{cousinweb}
\textbf{\hi{महफ़िल}} (\emph{mahfil}) is usually traced to Arabic \emph{maḥfil}, \textquotedblleft{}a gathering\textquotedblright{}, through Persian.
\end{cousinweb}

\begin{grammarlens}[title={What you've built}]
One more word for describing what you lived and felt.
\end{grammarlens}

\subsection*{Your turn}
\begin{itemize}
\raggedright
  \item \emph{Say it:} \emph{mahfil}
  \item \emph{Say it:} \emph{mahfil}, once more
  \item \emph{Say it:} \emph{Śām ko ghar par saṅgīt kī mahfil thī.} — In the evening there was a musical gathering at home.
  \item \emph{Recall:} say the Hindi for tempo, then the Hindi for practice (of music), then say \emph{mahfil} again
\end{itemize}

\begin{tcolorbox}[breakable,colback=teal!4,colframe=teal!35!black,title={Before you move on}]
What does \hi{महफ़िल} mean? (\textquotedblleft{}A gathering (for music or poetry)\textquotedblright{}.) Say it once more.
\end{tcolorbox}
\section[\texorpdfstring{\hi{नचाना} (nacānā)}{नचाना (nacānā)}]{\hi{नचाना} (nacānā) — to make dance}

\label{lesson:HI-C532-nacana}



\begin{quote}
Before the new one: say the Hindi for practice (of music), then the Hindi for a gathering (for music or poetry).
\end{quote}

\subsection*{You'll want to know: \hi{नचाना}}
\textbf{\hi{नचाना}} — \emph{nacānā} — \textquotedblleft{}to make dance\textquotedblright{}.

\begin{quote}
\hi{ढोल} \hi{ने} \hi{सब} \hi{लोगों} \hi{को} \hi{नचाया।}
\end{quote}

\begin{quote}
\emph{Ḍhol ne sab logoṁ ko nacāyā.}
\end{quote}

\begin{quote}
The drum got everyone dancing.
\end{quote}

Past, with \hi{ने} (he or she …): \hi{उसने} \hi{बच्चों} \hi{को} \hi{नचाया}, \emph{usne baccoṁ ko nacāyā}.

\begin{cousinweb}
\textbf{\hi{नचाना}} (\emph{nacānā}) is the causative of \emph{nācnā}, \textquotedblleft{}to dance\textquotedblright{}, which comes from Sanskrit \emph{nṛtyati}, \textquotedblleft{}dances\textquotedblright{} (CDIAL 7583).
\end{cousinweb}

\begin{grammarlens}[title={What you've built}]
One more word for describing what you lived and felt.
\end{grammarlens}

\subsection*{Your turn}
\begin{itemize}
\raggedright
  \item \emph{Say it:} \emph{nacānā}
  \item \emph{Say it:} \emph{nacānā}, once more
  \item \emph{Say it:} \emph{Ḍhol ne sab logoṁ ko nacāyā.} — The drum got everyone dancing.
  \item \emph{Recall:} say the Hindi for practice (of music), then the Hindi for a gathering (for music or poetry), then say \emph{nacānā} again
\end{itemize}

\begin{tcolorbox}[breakable,colback=teal!4,colframe=teal!35!black,title={Before you move on}]
What does \hi{नचाना} mean? (\textquotedblleft{}To make dance\textquotedblright{}.) Say it once more.
\end{tcolorbox}
